\section{Conclusion}
\label{sec:conclusion}

We have conducted the first systematic study of inference-time self-critique for LLM story writing, measuring improvement across six quality dimensions, three iteration depths, and two strategy families (\selfrefine vs.\ \bestofn).
Our central finding is that a single \critiquerevise iteration produces a large, statistically significant improvement in story quality (80\% win rate, Cohen's $d = 2.02$), with the strongest gains in creativity and character development.
Additional iterations offer diminishing returns, and self-selection without structured critique is largely ineffective.

These results carry a clear practical message: applications using LLMs for creative writing should implement one self-critique pass by default.
This costs only 3$\times$ the base generation cost but produces dramatically better output.
More iterations are not cost-effective.

Our findings also suggest a broader insight about LLM creative capacity.
The model \emph{knows} what would make a better story---more original premises, deeper characters, stronger emotional beats---but does not produce it in first drafts.
Structured self-reflection provides a mechanism to surface and act on this latent knowledge.
The challenge for future work is to determine whether this insight extends to longer narratives, other creative domains, and weaker models, and whether training-time integration of self-critique can produce permanent improvements in creative quality.
